\documentclass[11pt,twoside,a4paper]{article}

\usepackage[utf8]{inputenc}
\usepackage[T1]{fontenc}
\usepackage[ngerman]{babel}

\usepackage{amsmath,amssymb,amsfonts,amsthm}
\usepackage{booktabs}
\usepackage{geometry}
\usepackage{array}
\usepackage{cite}
\usepackage{hyperref}

\geometry{
	tmargin=2.5cm,
	bmargin=2.5cm,
	lmargin=2.5cm,
	rmargin=2.5cm
}

\title{\textbf{Exakter isotroper Pullback des $E_8$-Wurzelsystems und geometrische Regularisierung der 3D-Navier-Stokes-Gleichungen im EABC-Framework}}
\author{\textbf{EABC-Audit-Konsortium} \\ \textit{im Auftrag der Forschungsgruppe}}
\date{18. August 2026}

\newtheorem{theorem}{Theorem}[section]
\newtheorem{lemma}[theorem]{Lemma}
\newtheorem{corollary}[theorem]{Corollary}
\newtheorem{definition}[theorem]{Definition}
\newtheorem{remark}[theorem]{Remark}

\begin{document}
	
	\maketitle
	
	\begin{abstract}
		In dieser Arbeit präsentieren wir eine rigorose, geometrische Regularitätsroute für die inkompressiblen dreidimensionalen Navier-Stokes-Gleichungen im Rahmen der EABC-Symmetrieleiter, speziell auf den Ebenen 9 und 11. Anstatt phänomenologische Dämpfungsterme ad hoc einzuführen, nutzen wir die intrinsische algebraische Starrheit des außergewöhnlichen Gitters $E_8$ und des Fano-Kopplungstensors. Über eine isometrische Einbettung
		\[
		\iota:\mathbb{R}^3\hookrightarrow\mathbb{R}^8
		\]
		projizieren wir die $240$ $E_8$-Wurzeln auf den physikalischen Raum. Das sphärische 7-Design des Gitters garantiert eine exakt isotrope Identität des vierten Moments, woraus über Plancherel das quartische Symbol
		\[
		\mathcal{Q}_{E_8}[\xi]
		=
		36\|D^2\xi\|_2^2
		\]
		folgt. In drei Raumdimensionen führt eine uniforme $H^2$-Kontrolle über die Sobolev-Morrey-Kette
		\[
		H^2(\mathbb{R}^3)
		\hookrightarrow
		W^{1,6}(\mathbb{R}^3)
		\hookrightarrow
		C^{0,1/2}(\mathbb{R}^3)
		\]
		zu einer Hölder-Kontrolle des Richtungsfeldes. Diese analytische Route ist mit dem klassischen Programm der Vortizitätsrichtungs-Kohärenz kompatibel. Zusätzlich werden lineare biharmonische und stochastisch erzwungene Testmodelle sowie ein numerisches CLM-Regularisierungsaudit untersucht. Schließlich betrachten wir die Weyl-invariante Starrheit des Deformationsraums
		\[
		\operatorname{Sym}_0(8),
		\qquad
		\dim\operatorname{Sym}_0(8)=35.
		\]
		Bewiesene algebraische Aussagen werden strikt von weiterhin offenen physikalisch-dynamischen PDE-Fragen getrennt.
	\end{abstract}
	
	\section{Einleitung und die Skalierungsbarriere}
	
	Die Frage nach globaler Glattheit und Existenz von Lösungen der inkompressiblen dreidimensionalen Navier-Stokes-Gleichungen gehört zu den bedeutendsten ungelösten Problemen der mathematischen Physik.
	
	Die Vortizitätsgleichung lautet
	\begin{equation}
		\partial_t\omega+(u\cdot\nabla)\omega
		=
		(\omega\cdot\nabla)u+\nu\Delta\omega,
	\end{equation}
	wobei
	\[
	\omega=\nabla\times u
	\]
	die Vortizität und $\nu>0$ die kinematische Viskosität bezeichnet.
	
	Für die Enstrophie gilt
	\begin{equation}
		\frac{1}{2}\frac{d}{dt}\|\omega\|_2^2
		+
		\nu\|\nabla\omega\|_2^2
		=
		\int_{\mathbb{R}^3}
		\omega\cdot S\omega\,dx,
	\end{equation}
	mit
	\[
	S=\frac{1}{2}
	\left(
	\nabla u+(\nabla u)^T
	\right).
	\]
	
	Unter der Navier-Stokes-Skalierung
	\[
	u_\lambda(x,t)
	=
	\lambda u(\lambda x,\lambda^2t)
	\]
	gilt
	\[
	\omega_\lambda(x,t)
	=
	\lambda^2\omega(\lambda x,\lambda^2t).
	\]
	
	Daraus folgen
	\[
	\|\omega_\lambda\|_2^2
	=
	\lambda\|\omega\|_2^2,
	\]
	während sowohl der Wirbelstreckungsterm als auch die viskose Enstrophiedissipation wie
	\[
	\lambda^3
	\]
	skalieren.
	
	Daher kann ein zusätzlicher rein amplitudenbasierter Term proportional zu
	\[
	\|\omega\|_2^2
	\]
	die kritische UV-Skalierung nicht allein kontrollieren.
	
	Das EABC-Framework untersucht deshalb eine geometrische Route über die normierte Wirbelrichtung
	\begin{equation}
		\xi(x,t)
		=
		\frac{\omega(x,t)}{|\omega(x,t)|},
		\qquad
		\omega(x,t)\neq0.
	\end{equation}
	
	Diese Strategie steht in Beziehung zum klassischen Regularitätsprogramm von Constantin und Fefferman~\cite{constantin1993}, in dem geometrische Kohärenz der Vortizitätsrichtung eine zentrale Rolle spielt.
	
	\section{Der Fano-Winkelmesser als Richtungsdetektor}
	
	Auf einem ausgewählten quaternionischen Unterraum
	\[
	\mathbb{H}_{\rm Im}\cong\mathbb{R}^3
	\]
	reduziert sich die oktonionische 3-Form auf
	\begin{equation}
		\phi_{ijk}
		=
		\sigma\epsilon_{ijk},
		\qquad
		\sigma\in\{-1,+1\}.
	\end{equation}
	
	Für zwei Vektoren
	\[
	\xi,\eta\in\mathbb{R}^3
	\]
	definieren wir
	\begin{equation}
		V_k
		=
		\sum_{i,j=1}^3
		\phi_{ijk}\xi_i\eta_j.
	\end{equation}
	
	Damit gilt
	\begin{equation}
		V
		=
		\sigma(\xi\times\eta).
	\end{equation}
	
	Folglich ist
	\begin{equation}
		\|V\|_2
		=
		\|\xi\times\eta\|_2.
	\end{equation}
	
	Sind $\xi$ und $\eta$ normiert, so folgt
	\begin{equation}
		\boxed{
			\|\phi(\xi,\eta,\cdot)\|
			=
			|\sin\theta|
		}
	\end{equation}
	für den Winkel $\theta$ zwischen $\xi$ und $\eta$.
	
	Diese Normidentität ist unabhängig vom Orientierungszeichen $\sigma$.
	
	\section{Der isometrische Pullback des $E_8$-Wurzelsystems}
	
	Sei
	\[
	R(E_8)\subset\mathbb{R}^8
	\]
	das Wurzelsystem von $E_8$ mit $240$ Wurzeln und
	\[
	|r|^2=2
	\qquad
	\text{für alle }r\in R(E_8).
	\]
	
	Die normierten $E_8$-Wurzeln bilden ein \textbf{sphärisches 7-Design}. Daher gilt für jedes
	\[
	k\in\mathbb{R}^8
	\]
	die exakte Identität
	\begin{equation}
		\boxed{
			\sum_{r\in R(E_8)}
			(r\cdot k)^4
			=
			36|k|^4.
		}
	\end{equation}
	
	Nun sei
	\begin{equation}
		\iota:
		\mathbb{R}^3
		\hookrightarrow
		\mathbb{R}^8
	\end{equation}
	eine isometrische Einbettung mit
	\begin{equation}
		\iota^T\iota
		=
		I_3.
	\end{equation}
	
	Wir definieren
	\begin{equation}
		q_r
		=
		\iota^Tr
		\in
		\mathbb{R}^3.
	\end{equation}
	
	Für
	\[
	\kappa\in\mathbb{R}^3
	\]
	gilt
	\begin{align}
		\sum_{r\in R(E_8)}
		(q_r\cdot\kappa)^4
		&=
		\sum_{r\in R(E_8)}
		(r\cdot\iota\kappa)^4
		\\
		&=
		36|\iota\kappa|^4
		\\
		&=
		36|\kappa|^4.
	\end{align}
	
	Damit folgt die dimensionskonsistente Pullback-Identität
	\begin{equation}
		\boxed{
			\sum_{r\in R(E_8)}
			(q_r\cdot\kappa)^4
			=
			36|\kappa|^4.
		}
	\end{equation}
	
	\section{Die $H^2$-Kovarianzsteuerung und die Plancherel-Identität}
	
	Für ein hinreichend reguläres Richtungsfeld
	\[
	\xi:
	\mathbb{R}^3
	\to
	\mathbb{R}^3
	\]
	definieren wir
	\begin{equation}
		\mathcal{Q}_{E_8}[\xi]
		=
		\sum_{r\in R(E_8)}
		\left\|
		(q_r\cdot\nabla)^2\xi
		\right\|_{L^2(\mathbb{R}^3)}^2.
	\end{equation}
	
	Im Fourier-Raum gilt
	\begin{align}
		\mathcal{Q}_{E_8}[\xi]
		&=
		\sum_{r\in R(E_8)}
		\int_{\mathbb{R}^3}
		(q_r\cdot\kappa)^4
		|\widehat{\xi}(\kappa)|^2
		\,d\kappa
		\\
		&=
		36
		\int_{\mathbb{R}^3}
		|\kappa|^4
		|\widehat{\xi}(\kappa)|^2
		\,d\kappa.
	\end{align}
	
	Mit der üblichen Fourier-Normierung folgt
	\begin{equation}
		\boxed{
			\mathcal{Q}_{E_8}[\xi]
			=
			36
			\|D^2\xi\|_{L^2(\mathbb{R}^3)}^2.
		}
	\end{equation}
	
	Diese Identität liefert ein exakt isotropes Symbol vierter Ordnung.
	
	\section{Der Morrey-Sicherheitskorridor}
	
	In drei Raumdimensionen gilt
	\begin{equation}
		H^2(\mathbb{R}^3)
		\hookrightarrow
		W^{1,6}(\mathbb{R}^3).
	\end{equation}
	
	Weiter gilt in geeigneten beschränkten oder lokalisierten Gebieten die Morrey-Einbettung
	\begin{equation}
		W^{1,6}
		\hookrightarrow
		C^{0,1/2}.
	\end{equation}
	
	Damit ergibt sich unter geeigneten Lokalisierungs- und Normannahmen die konditionale Kette
	\begin{equation}
		\boxed{
			\mathcal{Q}_{E_8}[\xi]
			\le C
			\quad\Longrightarrow\quad
			H^2\text{-Kontrolle}
			\quad\Longrightarrow\quad
			W^{1,6}\text{-Kontrolle}
			\quad\Longrightarrow\quad
			C^{0,1/2}\text{-Kontrolle}.
		}
	\end{equation}
	
	Dabei ist zu beachten, dass
	\[
	\xi
	=
	\frac{\omega}{|\omega|}
	\]
	auf der Nullvortizitätsmenge nicht definiert ist. Eine Navier-Stokes-Anwendung muss deshalb auf geeigneten High-Vorticity-Regionen oder über eine mathematisch kontrollierte Fortsetzung formuliert werden.
	
	\section{Quantitative Simulations- und Audit-Ergebnisse}
	
	\subsection{1D-Simulationsaudit des biharmonischen Flusses}
	
	Als lineares Testmodell betrachten wir
	\begin{equation}
		d\eta
		=
		-36\gamma
		\partial_x^4\eta
		\,dt
		+
		\sigma\,dW
	\end{equation}
	auf einem periodischen eindimensionalen Gebiet.
	
	Für jede Fouriermode gilt
	\begin{equation}
		d\widehat{\eta}_k
		=
		-36\gamma k^4
		\widehat{\eta}_k
		\,dt
		+
		\sigma_k\,dW_k.
	\end{equation}
	
	Ohne Rauschen besitzt die Mode die exakte Lösung
	\begin{equation}
		\widehat{\eta}_k(t)
		=
		e^{-36\gamma k^4t}
		\widehat{\eta}_k(0).
	\end{equation}
	
	Die Relaxationszeit ist somit
	\begin{equation}
		\boxed{
			\tau_k
			=
			\frac{1}{36\gamma k^4}.
		}
	\end{equation}
	
	Für den numerischen Audit wurden die Parameter
	\[
	\gamma=0{,}5,
	\qquad
	\sigma=0{,}05,
	\qquad
	N=256,
	\qquad
	T=0{,}5
	\]
	verwendet.
	
	Der gemessene spätere spektrale Log-Log-Exponent betrug
	\begin{equation}
		\alpha_{\rm emp}
		=
		-3{,}9571,
	\end{equation}
	gegenüber dem theoretischen Wert
	\[
	-4.
	\]
	
	\begin{table}[htbp]
		\centering
		\caption{Biharmonische Relaxationszeiten für $\gamma=0{,}5$.}
		\label{tab:relaxation_times}
		\begin{tabular}{crr}
			\toprule
			$k$ &
			$36\gamma k^4$ &
			$\tau_k$
			\\
			\midrule
			1  & 18 & $5{,}56\times10^{-2}$ \\
			2  & 288 & $3{,}47\times10^{-3}$ \\
			4  & 4608 & $2{,}17\times10^{-4}$ \\
			8  & 73728 & $1{,}36\times10^{-5}$ \\
			16 & $1{,}18\times10^6$ & $8{,}48\times10^{-7}$ \\
			\bottomrule
		\end{tabular}
	\end{table}
	
	\subsection{Regularisierungs-Audit des Constantin-Lax-Majda-Modells}
	
	Als nichtlineares eindimensionales Testmodell betrachten wir
	\begin{equation}
		\partial_t\omega
		=
		\omega H(\omega)
		-
		36\gamma
		\partial_x^4\omega,
	\end{equation}
	wobei $H$ die Hilbert-Transformation bezeichnet.
	
	Für die konkret untersuchte numerische Konfiguration wurde per Bisektion ein Schwellenwert
	\begin{equation}
		\gamma_{\rm crit}
		\approx
		0{,}004199
	\end{equation}
	ermittelt.
	
	Dieser Wert ist
	\begin{center}
		\textbf{modell-, daten-, diskretisierungs- und solverabhängig}
	\end{center}
	und stellt keine universelle analytische Konstante dar.
	
	Aus diesem numerischen CLM-Audit folgt insbesondere kein Beweis der globalen 3D-Navier-Stokes-Regularität.
	
	\section{Das $E_8$-Weyl-Starrheitsaudit}
	
	Der euklidische Shape-Raum sei
	\begin{equation}
		\mathbf{S}
		=
		\operatorname{Sym}_0(8),
	\end{equation}
	also der Raum der reellen symmetrischen spurfreien $8\times8$-Matrizen.
	
	Es gilt
	\begin{equation}
		\dim\mathbf{S}
		=
		\frac{8\cdot9}{2}-1
		=
		35.
	\end{equation}
	
	Für Weyl-Reflexionen
	\begin{equation}
		s_\alpha
		=
		I_8
		-
		\alpha\alpha^T,
		\qquad
		|\alpha|^2=2,
	\end{equation}
	wirkt die induzierte Darstellung durch
	\begin{equation}
		M
		\mapsto
		s_\alpha
		M
		s_\alpha^T.
	\end{equation}
	
	Das vorbereitete exakte Audit untersucht die beiden Bedingungen
	\begin{equation}
		\operatorname{Fix}_{W(E_8)}
		\left(
		\operatorname{Sym}_0(8)
		\right)
		=
		\{0\}
	\end{equation}
	und
	\begin{equation}
		\dim
		\left(
		\operatorname{Sym}^2
		\left(
		\operatorname{Sym}_0(8)^*
		\right)^{W(E_8)}
		\right)
		=
		1.
	\end{equation}
	
	Für das explizite Fisher-Umegaki-Shape-Modell wurde unabhängig davon die quadratische Hesseform
	\begin{equation}
		\boxed{
			H_{\rm Fisher}
			=
			\frac{1}{2}
			I_{35}
		}
	\end{equation}
	bestätigt.
	
	Damit ist dieses spezielle quadratische Modell am symmetrischen Punkt isotrop.
	
	Höhergradige Weyl-Invarianten werden dadurch nicht ausgeschlossen.
	
	\section{Schlussfolgerung und Claim-Wall}
	
	Der aktuelle epistemische Status lautet:
	
	\begin{table}[htbp]
		\centering
		\caption{Status des Fano-$E_8$-Navier-Stokes-Zweigs.}
		\label{tab:claim_wall}
		\begin{tabular}{
				>{\raggedright\arraybackslash}p{0.16\textwidth}
				>{\raggedright\arraybackslash}p{0.53\textwidth}
				>{\raggedright\arraybackslash}p{0.21\textwidth}
			}
			\toprule
			Gate &
			Relation &
			Status
			\\
			\midrule
			
			$G_{\rm NS}^{(9a)}$
			&
			$E_8$-Pol bei $s=8$
			&
			\textbf{EXACT/STANDARD}
			\\
			
			$G_{\rm NS}^{(9b)}$
			&
			$\displaystyle
			\sum_r(r\cdot k)^4
			=
			36|k|^4$
			&
			\textbf{EXACT PASS}
			\\
			
			$G_{\rm NS}^{(9c)}$
			&
			$\displaystyle
			\sum_r(q_r\cdot\kappa)^4
			=
			36|\kappa|^4$
			&
			\textbf{EXACT PASS}
			\\
			
			$G_{\rm NS}^{(9e)}$
			&
			$\displaystyle
			\mathcal{Q}_{E_8}[\xi]
			=
			36\|D^2\xi\|_2^2$
			&
			\textbf{EXACT PASS}
			\\
			
			$G_{\rm NS}^{(9f)}$
			&
			$H^2\to W^{1,6}\to C^{0,1/2}$
			&
			\textbf{STANDARD}
			\\
			
			$G_{\rm NS}^{(9g)}$
			&
			Vortizitätsrichtungs-Regularitätsbrücke
			&
			\textbf{CONDITIONAL}
			\\
			
			$G_{\rm NS}^{(9d)}$
			&
			Dynamische Herleitung und a-priori-Kontrolle von $\gamma>0$
			&
			\textbf{OPEN}
			\\
			
			---
			&
			Globale 3D-Navier-Stokes-Regularität
			&
			\textbf{OPEN}
			\\
			
			\bottomrule
		\end{tabular}
	\end{table}
	
	Die algebraischen $E_8$-Identitäten und die Plancherel-Reduktion sind exakt. Offen bleiben insbesondere die dynamische Herleitung einer positiven PDE-Kopplung $\gamma$, ihre uniforme Kontrolle entlang der tatsächlichen Strömungsentwicklung sowie die vollständige Übertragung auf die dreidimensionale Navier-Stokes-Dynamik.
	
	\appendix
	
	\section{Anhang: Formale Lean-4-Spezifikation}
	\label{sec:appendix_lean_formal}
	
	Zur formalen Dokumentation geben wir im Folgenden ein ASCII-kompatibles Lean-4-Skelett an. Der Block ist bewusst frei von Unicode-Sonderzeichen gehalten, damit er mit pdfLaTeX innerhalb einer \texttt{verbatim}-Umgebung kompiliert werden kann.
	
	\begin{verbatim}
		import Mathlib.LinearAlgebra.CliffordAlgebra.Basic
		import Mathlib.LinearAlgebra.QuadraticForm.Basic
		
		/-!
		# EABC Clifford vacuum and Fano detector
		
		Formalization target for:
		1. the double-sheet Clifford structure,
		2. the Fano angle identity on a 3D quaternionic sector.
		
		This is a Lean 4 target skeleton.
		Exact Mathlib API compatibility must be checked separately.
		-/
		
		open CliffordAlgebra
		
		section DoubleSheetClifford
		
		variable (R : Type*) [CommRing R] [Algebra Real R]
		
		structure DoubleSheetSpace where
		gamma : Real
		j_switch : Real
		
		def doubleSheetQForm :
		QuadraticForm Real DoubleSheetSpace :=
		QuadraticForm.mk
		(fun v => v.gamma ^ 2 - v.j_switch ^ 2)
		
		def CliffordDoubleSheet :=
		CliffordAlgebra (doubleSheetQForm Real)
		
		variable (v_gamma v_j : DoubleSheetSpace)
		
		def is_orthogonal_generators : Prop :=
		QuadraticForm.IsOrthogonal
		(doubleSheetQForm Real) v_gamma v_j
		
		def Gamma_op :
		CliffordDoubleSheet Real :=
		iota (doubleSheetQForm Real) v_gamma
		
		def J_op :
		CliffordDoubleSheet Real :=
		iota (doubleSheetQForm Real) v_j
		
		theorem anti_commute_generators
		(h_orth :
		is_orthogonal_generators Real v_gamma v_j) :
		Gamma_op Real v_gamma * J_op Real v_j
		+ J_op Real v_j * Gamma_op Real v_gamma
		= 0 := by
		sorry
		
		def K_op :
		CliffordDoubleSheet Real :=
		Gamma_op Real v_gamma * J_op Real v_j
		
		theorem complex_structure_squared_neg_one
		(h_orth :
		is_orthogonal_generators Real v_gamma v_j)
		(h_norm_gamma :
		(doubleSheetQForm Real) v_gamma = -1)
		(h_norm_j :
		(doubleSheetQForm Real) v_j = 1) :
		K_op Real v_gamma v_j * K_op Real v_gamma v_j
		= -1 := by
		sorry
		
		end DoubleSheetClifford
		
		
		section FanoAngleDetector
		
		structure ImH where
		x : Real
		y : Real
		z : Real
		
		def fano_det (xi eta : ImH) : ImH :=
		{
			x := xi.y * eta.z - xi.z * eta.y
			y := xi.z * eta.x - xi.x * eta.z
			z := xi.x * eta.y - xi.y * eta.x
		}
		
		def norm_L2 (v : ImH) : Real :=
		Real.sqrt (v.x^2 + v.y^2 + v.z^2)
		
		def dot3 (xi eta : ImH) : Real :=
		xi.x * eta.x
		+ xi.y * eta.y
		+ xi.z * eta.z
		
		theorem fano_angle_identity
		(xi eta : ImH)
		(hxi : norm_L2 xi = 1)
		(heta : norm_L2 eta = 1) :
		norm_L2 (fano_det xi eta)
		=
		Real.sqrt
		(1 - (dot3 xi eta)^2) := by
		sorry
		
		end FanoAngleDetector
	\end{verbatim}
	
	\begin{thebibliography}{9}
		
		\bibitem{constantin1993}
		P. Constantin and C. Fefferman,
		\emph{Direction of vorticity and the problem of global regularity for the Navier-Stokes equations},
		Indiana University Mathematics Journal,
		42(3), 775--789 (1993).
		
		\bibitem{chemtov2022}
		M. Chemtov and S. Karigiannis,
		\emph{Canonical forms of skew-symmetric operators in $G_2$ geometry},
		Journal of Geometry and Physics,
		178, 104553 (2022).
		
	\end{thebibliography}
	
\end{document}